
%\documentclass[gray,trans]{beamer} % com o mod transpar�ncia
%\documentclass[gray]{beamer}
\documentclass[trans,note]{beamer}

\usetheme{Madrid}

%\usetheme{Marburg} % Template ideial para SBMF
%\usetheme{CambridgeUS}
 %\usetheme{Boadilla}
%\usetheme{default}


%\usecolortheme{structure} 
%\usebeamercolor{structure}

% \usefonttheme[onlylarge]{structurebold}
% \usefonttheme{default}
% \setbeamerfont*{frametitle}{size=\normalsize,series=\bfseries}
 % \setbeamertemplate{navigation symbols}{}

% \usepackage[latin1]{inputenc}
% \usepackage[brazil]{babel}
\usepackage[english]{babel}
\usepackage[latin1]{inputenc}

\usepackage{graphicx,url}
\usepackage{fancyhdr}

\usepackage{tikz}
\usetikzlibrary{arrows}
 \tikzstyle{block}=[draw opacity=0.3,line width=1.0cm]

% Class options include: notes, notesonly, handout, trans,
% hidesubsections, shadesubsections,
% inrow, blue, red, grey, brown

% Theme for beamer presentation.
%\usepackage{beamerthemeclassic}
% \usepackage{beamerthemebars}
% \usepackage{beamerthemesplit}
% \usepackage{beamerthemetree}
% \usepackage{beamerthemelined}

\title{Formal Modelling of a Microcontroller Instruction Set in B}

%\subtitle{Qualifica��o de Mestrado} % Enter your title betweencurlybraces
\author{Val�rio Medeiros Jr, David D�harbe}
\institute{Federal University of Rio Grande do Norte} % Enter your institute name between curly braces
\date{20 August 2009} % Enter the date or \today between curly braces

\begin{document}

%% Keywords for the B method
\newcommand{\MACHINE}{\operatorname{\mathbf{MACHINE}}}
\newcommand{\REFINEMENT}{\operatorname{\mathbf{REFINEMENT}}}
\newcommand{\IMPLEMENTATION}{\operatorname{\mathbf{IMPLEMENTATION}}}
\newcommand{\REFINES}{\operatorname{\mathbf{REFINES}}}
\newcommand{\SEES}{\operatorname{\mathbf{SEES}}}
\newcommand{\INCLUDES}{\operatorname{\mathbf{INCLUDES}}}
\newcommand{\IMPORTS}{\operatorname{\mathbf{IMPORTS}}}
\newcommand{\SETS}{\operatorname{\mathbf{SETS}}}
\newcommand{\CONSTANTS}{\operatorname{\mathbf{CONSTANTS}}}
\newcommand{\PROPERTIES}{\operatorname{\mathbf{PROPERTIES}}}
\newcommand{\CONCRETE}{\operatorname{\mathbf{CONCRETE}}}
\newcommand{\VARIABLES}{\operatorname{\mathbf{VARIABLES}}}
\newcommand{\ASSERTIONS}{\operatorname{\mathbf{ASSERTIONS}}}
\newcommand{\CONCRETEVARIABLES}{\operatorname{\mathbf{CONCRETE\_VARIABLES}}}
\newcommand{\DEFINITIONS}{\operatorname{\mathbf{DEFINITIONS}}}
\newcommand{\VAR}{\operatorname{\mathbf{VAR}}}
\newcommand{\IN}{\operatorname{\mathbf{IN}}}
\newcommand{\INVARIANT}{\operatorname{\mathbf{INVARIANT}}}
\newcommand{\INITIALISATION}{\operatorname{\mathbf{INITIALISATION}}}
\newcommand{\OPERATIONS}{\operatorname{\mathbf{OPERATIONS}}}
\newcommand{\BEGIN}{\operatorname{\mathbf{BEGIN}}}
\newcommand{\END}{\operatorname{\mathbf{END}}}
\newcommand{\PRE}{\operatorname{\mathbf{PRE}}}
\newcommand{\IF}{\operatorname{\mathbf{IF}}}
\newcommand{\THEN}{\operatorname{\mathbf{THEN}}}
\newcommand{\ELSE}{\operatorname{\mathbf{ELSE}}}
\newcommand{\ELSIF}{\operatorname{\mathbf{ELSIF}}}
\newcommand{\ANY}{\operatorname{\mathbf{ANY}}}
\newcommand{\WHERE}{\operatorname{\mathbf{WHERE}}}
\newcommand{\CASE}{\operatorname{\mathbf{CASE}}}
\newcommand{\OF}{\operatorname{\mathbf{OF}}}
\newcommand{\EITHER}{\operatorname{\mathbf{EITHER}}}
\newcommand{\AND}{\operatorname{\mathbf{AND}}}
\newcommand{\OR}{\operatorname{\mathbf{OR}}}
\newcommand{\NOT}{\operatorname{\mathbf{NOT}}}
\newcommand{\WHILE}{\operatorname{\mathbf{WHILE}}}
\newcommand{\DO}{\operatorname{\mathbf{DO}}}
\newcommand{\VARIANT}{\operatorname{\mathbf{VARIANT}}}
\newcommand{\FALSE}{\operatorname{\mathbf{FALSE}}}
\newcommand{\TRUE}{\operatorname{\mathbf{TRUE}}}

%% Commonly used math entities
\newcommand{\pow}{\operatorname{\mathbb{P}}}
\newcommand{\nat}{\operatorname{\mathbb{N}}}
\newcommand{\pfun}{\operatorname{\rightarrow\mkern-22mu+}}
\newcommand{\fset}{\operatorname{\mathbb{F}}}
\newcommand{\dom}{\operatorname{\mbox{dom}}}
\newcommand{\ran}{\operatorname{\mbox{ran}}}
\newcommand{\natone}{\operatorname{\mathbb{N}_1}}
\newcommand{\integer}{\operatorname{\mathbb{Z}}}
\newcommand{\fun}{\operatorname{\rightarrow}}
\newcommand{\domr}{\operatorname{\triangleleft}}
\newcommand{\seq}{\operatorname{\mathbf{seq1}}}
\newcommand{\ovr}{\operatorname{\oplus}}
\newcommand{\BOOL}{\operatorname{\mathbf{BOOL}}}
\newcommand{\pred}{\operatorname{\mathbf{pred}}}
\newcommand{\Bsucc}{\operatorname{\mathbf{succ}}}


%\usepackage{supertabular}

% DEFINITION DES CARACTERES MATHEMATIQUES B
%------------------------------------------
\def\@setmcodes#1#2#3{{\count0=#1 \count1=#3
	\loop \global\mathcode\count0=\count1 \ifnum \count0<#2
	\advance\count0 by1 \advance\count1 by1 \repeat}}

%\@setmcodes{`A}{`Z}{"7441}
%\@setmcodes{`a}{`z}{"7461}

\mathcode`\;="8000 % Makes ; active in math mode
{\catcode`\;=\active \gdef;{\semicolon\;}}
\mathchardef\semicolon="003B
%    Nominal distance from top of paper to top of page
% \topmargin 0 pt
% \textheight 53\baselineskip
% 
% %   Left margin on odd-numbered pages
% \oddsidemargin  0.15 in
% %   Left margin on even-numbered pages
% \evensidemargin 0.35 in
% %   Width of marginal notes.
% \marginparwidth 1 in
% %   Note that \oddsidemargin = \evensidemargin
% \oddsidemargin 0.25 in
% \evensidemargin 0.25 in
% \marginparwidth 0.75 in
% \textwidth 5.875 in % Width of text line.
% 
% \setlength{\parindent}{0pt}
% \setlength{\parskip}{0ex}

% DEFINITION DES FONTS
%---------------------
% The AMS extra symbol fonts are loaded.
% Note: sometimes called euxm10
\font\msx=msam10
% Note: sometimes called euym10
\font\msy=msbm10

\newfam\msxfam \textfont\msxfam=\msx
\newfam\msyfam \textfont\msyfam=\msy

\def\famletter#1{\ifcase #1 0\or 1\or 2\or 3\or 4\or 5\or 6\or 7\or
	8\or 9\or A\or B\or C\or D\or E\or F\fi}

\edef\fx{\famletter\msxfam}
\edef\fy{\famletter\msyfam}

\def\bbold{\fam\msyfam \msy}

% SYMBOLES B
%-----------
% makes a quoted expression in mathematical text
\def\token#1{\hbox{`$#1$'}}
% used for error messages in Z specs
\def\report#1{\hbox{`{\tt #1}'}}

% \@myop makes an operator, with a strut to defeat TeX's vertical adjustment.
\def\@myop#1{\mathop{\mathstrut{#1}}\nolimits}

% This underscore doesn't have the little kern --- you get an italic
% correction anyway in math mode.
\def\_{\leavevmode \vbox{\hrule width0.5em}}

% Save \q as \xq for quantifiers q.
\let\xforall=\forall
\let\xexists=\exists
\let\xlambda=\lambda
\let\xmu=\mu

% \p and \f make arrows with 1 and 2 crossings resp.
\def\p#1{\mathrel{\ooalign{\hfil$\mapstochar\mkern 5mu$\hfil\cr$#1$}}}
\def\f#1{\mathrel{\ooalign{\hfil
	$\mapstochar\mkern 3mu\mapstochar\mkern 5mu$\hfil\cr$#1$}}}

\let\mc=\mathchardef

\def	\pow		{\mbox{${\cal P}$}}
\def	\po1		{\mbox{${\cal P}_1$}}
\let	\cross		\times
\def	\lambda		{\@myop{\xlambda}}
\def	\lnot		{\neg\;}
\def	\land		{\mathrel{\wedge}}
\def	\lor		{\mathrel{\vee}}
\let	\implies	\Rightarrow
\let	\iff		\Leftrightarrow
\def	\forall		{\@myop{\xforall}}
\def	\exists		{\@myop{\xexists}}
\def	\semi		{\mathrel{\comp}}
\def	\ssemi		{\mathbin{\rm ;}}
\let	\ensembleVide	\emptyset
\let	\rel		\leftrightarrow
\def	\dom		{\@myop{\sf dom}}
\def	\ran		{\@myop{\sf ran}}
\def	\id		{\@myop{\sf id}}
\def	\comp		{\mathbin{\raise
			0.6ex\hbox{\oalign{\hfil$\scriptscriptstyle
			\rm o$\hfil\cr\hfil$\scriptscriptstyle\rm 9$\hfil}}}}
\def	\para		{\mbox{$\mid\mid$}}
\mc	\dres		"2\fx43
\mc	\rres		"2\fx42
\def	\ndres		{\mathbin{{\dres} \llap{$-$}}}
\def	\nrres		{\mathbin{{\rres}\llap{$-$}}}
\def	\lover		{\vartriangleleft{ \llap{$-\!\!\!\!-\!$}}}
\def	\rover		{\mathbin{{\rres}\llap{$\!-\!\!\!-$}}}
\let	\fun		\rightarrow
\def	\pfun		{\p\fun}
\def	\pinj		{\p\inj}
\mc	\inj		"3\fx1A
\def	\psurj		{\p\surj}
\mc	\surj		"3\fx10
\def	\bij		{\surj\!\!\!\!\!\!\!\inj}
\def	\nat		{\mbox{${\cal N}$}}
\def	\na1		{\mbox{${\cal N}_1$}}
\def	\num		{\mbox{${\cal Z}$}}
\def	\int		{\mbox{${\cal Z}$}}
\def	\rat		{\mbox{${\cal Q}$}}
\def	\div		{\mathbin{\rm /}}
\def	\mod		{\mathbin{\bf mod}}
\def	\upto		{\mathbin{\ldotp\ldotp}}
\def	\finset		{\mbox{${\cal F}$}}
\def	\finse1		{\mbox{${\cal F}_1$}}
\def	\ffun		{\f\fun}
\def	\finj		{\f\inj}
\def	\seq		{\@myop{\rm seq}}
\def	\cat		{\mathbin{\raise 0.8ex\hbox{$\mathchar"2\fx61$}}}
\def	\sep		{\hspace*{.05in}}

\setcounter{secnumdepth}{0}
\setcounter{tocdepth}{0}

%-------------------%
% Debut du document %
%-------------------%





\beamertemplatefootpagenumber % Mostra o n�mero de p�ginas

% Creates title page of slide show using above information
\begin{frame}
   \titlepage  
  \begin{center} 

  Project: B2ASM \\
  \includegraphics[width=.14\textwidth]{figuras/ufrn.png} \ \ \ \ \ \ \ \ \ \
  \ \ \ \ \ \ \ \ \ \ \ \ \ \ \ \ \ \ \ \ \ \ \ \ \ \ \ \  
   \includegraphics[width=.12\textwidth]{figuras/dimap.png}
  \end{center}
  
\end{frame}
\note{Slide Principal} % Add notes to yourself that will be displayed when
  % typeset with the notes or notesonly class options


%\section [Agenda]{Agenda}


% Creates table of contents slide incorporating
% all \section and \subsection commands
 \begin{frame}
  \frametitle{Agenda}
  \tableofcontents
 \end{frame}

% Mostra um slide com destaque a atual se��o, em cada mudan�a de se��o
\AtBeginSection[]
{
\begin{frame}
\frametitle{Summary}
\tableofcontents[currentsection]
\end{frame}
}

 \include{content/introduction}
 % * Explain the cover of B Method 
 % The solution is a framework proposed a paper [Dantas, 2009]
 % The component key to this aproach is the formal model
 % Then this paper gave u

 \include{content/metodoB}
 % 

 \include{content/Modelling}

 \include{content/Description}

 \include{content/Proofs}
 
 \include{content/Related_works}

  \include{content/conclusions}



% 
% 
% \begin{frame}
% \frametitle{Refer�ncias} % Insert frame title between curly braces
%  
% \begin{itemize}
% \item J. R. Abrial. ``The B Book: Assigning Programs
 % to Meanings''. 1. ed. United States of America:
% Cambridge University Press, 1996.
% 
% \item D. Bell, Ian Morrey e John Pugh, ``Software Engineering: A Programming
% Approach", segunda edi��o, Prentice Hall, Nova Iorque, 1992.
 % 
% \item DANTAS, B. P. et al. ``Applying the b method to take on the grand
% challenge of verified compilation''. In: SBMF. \textit{Brazilian Symposium on Formal Methods}.
% Salvador - BA, 2008.  
 % 
% \item DANTAS, B. P. et al, ``Proposta e Avalia��o de uma Abordagem de Desenvolvimento de Software
% Fidedigno por Constru��o com o M�todo B'', SEMISH, 2008, B�lem.
%  
% \end{itemize}
% 
 % \end{frame}
% \note{Referencias} 
% 
% 
% \begin{frame}
% \frametitle{Refer�ncias} % Insert frame title between curly braces
% \begin{itemize}
% \item Microchip. PIC16C432 - ``Data Sheet OTP 8-Bit CMOS MCU with LIN
 % Transceiver''. USA, 2002. 197p.
% \item D.E.C. Nicolosi, ``Microcontrolador 8051 detalhado'' ,Ed. Erica, S�o
% Paulo, 2005.
% \item B. William, ``Z80 microcomputer design projects'',Ed. Sams, 1980.
 % \item UK COMPUTING RESEARCH COMMITTEE.``The Verifying Compiler: A GrandChallenge for Computing Research''.
% In UK: Workshop on Grand Challenges for Computing Research, 2005. Dispon�vel
% em:http://www.nesc.ac.uk/esi/events/Grand Challenges/workshop02.html Acessado em:18 abr 2009.  
 %  
% \end{itemize}
%  
% \end{frame}
% \note{Referencias}
%  
  
   



\end{document}